% Glucose Dynamics Cycle Diagram
\begin{tikzpicture}[
  every node/.append style={font=\small},
  src/.style={ellipse, draw=orange!70, fill=orange!15, minimum width=2.2cm,
              minimum height=0.9cm, align=center, thick},
  proc/.style={rectangle, rounded corners=3pt, draw=green!60!black, fill=green!12,
               minimum width=2.2cm, minimum height=0.9cm, align=center, thick},
  cons/.style={rectangle, rounded corners=3pt, draw=red!60, fill=red!12,
               minimum width=2.2cm, minimum height=0.9cm, align=center, thick},
  eff/.style={rectangle, rounded corners=3pt, draw=blue!60, fill=blue!12,
              minimum width=2.2cm, minimum height=0.9cm, align=center, thick},
  arrow/.style={->, >=stealth, thick, shorten >=2pt, shorten <=2pt},
]

% 7 nodes arranged in a circle (radius ~3.2cm, starting from top, clockwise)
% Angles: 90, 38.6, -12.9, -64.3, -115.7, -167.1, 141.4
\def\R{3.2}

\node[src]  (bv)   at (90:\R)     {Blood\\[-1pt]Vessels};
\node[proc] (inj)  at (38.6:\R)   {Glucose\\[-1pt]Injection};
\node[proc] (diff) at (-12.9:\R)  {Spatial\\[-1pt]Diffusion};
\node[cons] (tum)  at (-64.3:\R)  {Tumor\\[-1pt]Consumption};
\node[cons] (star) at (-115.7:\R) {Immune\\[-1pt]Starvation};
\node[eff]  (grad) at (-167.1:\R) {Concentration\\[-1pt]Gradient};
\node[eff]  (chem) at (141.4:\R)  {Immune\\[-1pt]Chemotaxis};

% Cycle arrows
\draw[arrow, orange!70]       (bv)   -- (inj);
\draw[arrow, green!60!black]  (inj)  -- (diff);
\draw[arrow, green!60!black]  (diff) -- (tum);
\draw[arrow, red!60]          (tum)  -- (star);
\draw[arrow, red!60]          (star) -- (grad);
\draw[arrow, blue!60]         (grad) -- (chem);
\draw[arrow, blue!60]         (chem) -- (bv);

% Center label
\node[font=\small\bfseries, align=center, text=gray!70!black] at (0, 0) {Glucose\\Dynamics\\Cycle};

% Subtle annotation labels
\node[font=\footnotesize, orange!70!black, above=1pt] at (bv.north) {source};
\node[font=\footnotesize, red!60!black, below=1pt] at (tum.south) {Warburg effect};
\node[font=\footnotesize, blue!60!black, left=1pt] at (grad.west) {spatial signal};

\end{tikzpicture}
